\documentclass[skripsi]{unnes-ta_plus}

\unnesSetTitle{Judul Skripsi Anda}
\unnesSetDegree{Sarjana \dots}
\unnesSetAuthor{Nama Mahasiswa}
\unnesSetNIM{1234567890}
\unnesSetProdi{Teknik Informatika}
\unnesSetFakultas{Fakultas Matematika dan Ilmu Pengetahuan Alam}
\unnesSetCity{Semarang}
\unnesSetYear{2026}
\unnesSetPembimbing{Nama Pembimbing}{NIP Pembimbing}

\begin{document}
\unnesFrontMatter

% Bagian awal skripsi (panduan: sampul, halaman judul, persetujuan, pengesahan penguji, pernyataan, dst.)
\maketitleunnes
\unnesPersetujuanPembimbing

% Contoh: isi manual nama+NIP penguji
\unnesPengesahanPenguji
  {Nama Ketua Penguji / NIP}
  {Nama Sekretaris / NIP}
  {Nama Penguji 1 / NIP}
  {Nama Penguji 2 / NIP}

\unnesPernyataanKeaslian

\begin{unnesabstract}
% Abstrak spasi 1.0 (single). Isi komponen: judul, penulis, pembimbing, tujuan, metode, hasil/temuan.
Tuliskan abstrak di sini.
\end{unnesabstract}

\begin{unnesabstracten}
Write your English abstract here.
\end{unnesabstracten}

\tableofcontents
\clearpage

\unnesMainMatter

\chapter{Pendahuluan}
\chapter{Tinjauan Pustaka}
\chapter{Metode Penelitian}
\chapter{Hasil dan Pembahasan}
\chapter{Simpulan dan Saran}

\chapter*{Daftar Pustaka}
\addcontentsline{toc}{chapter}{Daftar Pustaka}

\end{document}
